%% ===========================================================================
%% Controlled-prototype results for the stage crossover.
%% Inserted as the first subsection of Section~\ref{sec:results}.
%%
%% DUPLICATION NOTE. paper/results_controlled.tex reports the same measurements against the
%% notation and equation labels of the Overleaf project's main_yd.tex (Yashar's rewrite:
%% \Dist_s, eq:direct/eq:state/eq:violation/eq:repair, RQ1--RQ4). This file is the counterpart
%% for the main.tex lineage that fair_trace_acm.tex belongs to. The two must be kept in sync
%% until one lineage wins; if main_yd.tex becomes the manuscript, delete this file and the
%% \input of it in fair_trace_acm.tex.
%%
%% All numbers are from the two-domain controlled prototype, seed 42, 24 scenarios x 4 repeats
%% per domain, no LLM judge. Reproduced from a local re-execution of the reference notebook,
%% which matched its recorded outputs exactly.
%% ===========================================================================

\subsection{Validating the Crossover on a Controlled Substrate}
\label{sec:scaff-results}

Before applying the decomposition of Section~\ref{sec:scaff} to an agent whose behaviour is
unknown, we validate that it recovers a dependence whose location we control. We construct two
controlled domains---movie and restaurant recommendation---in which the latent user profile,
the set of acceptable actions, item relevance, and memory support are all known by construction,
and in which a stage-specific attribute dependence is \emph{planted} at a known stage. The
planted dependence is an artificial unit test of the measurement instrument and carries no claim
about the preferences or treatment of any real group. Each domain contributes $24$ scenarios run
at $4$ repeats; every stage output is checked by a structural validity gate before any fairness
quantity is computed, so that an invalid pair is excluded rather than scored. Under this design a
stage sees the descriptor only when it is that scenario's planted stage, which means the causal
origin of every divergence is known and the decomposition has a ground truth to be judged
against.

Table~\ref{tab:scaff-coords} reports the mean per-stage coordinates. The pattern that motivates
the framework appears immediately at \stage{Rank}: it carries the largest natural divergence in
both domains ($0.396$ and $0.490$), and most of that divergence is inherited ($0.302$, $0.333$)
rather than direct ($0.094$, $0.156$). An audit reading $\Dnat_s$ alone would direct remediation
effort at the ranker. \stage{Retrieve} shows the opposite composition---direct exceeding
inherited in both domains---and is where the decomposition assigns responsibility.

\begin{table}[t]
  \caption{Mean per-stage coordinates on the controlled prototype ($24$ scenarios $\times\, 4$
    repeats per domain). \emph{noise} is the within-condition repeat distance and bounds what a
    direct effect must exceed to be interpretable. Diagnoses are the corpus-mean reading at
    $\tau = 0.10$ and summarize a distribution rather than delivering a per-conversation verdict.}
  \label{tab:scaff-coords}
  \begin{tabular}{llrrrrl}
    \toprule
    Domain & Stage & $\Dnat_s$ & $\DA_s$ & $\DZ_s$ & noise & diagnosis \\
    \midrule
    Movies      & \stage{Elicit}   & 0.056 & 0.056 & 0.000 & 0.000 & invariant \\
                & \stage{Retrieve} & 0.205 & 0.136 & 0.069 & 0.166 & direct \\
                & \stage{Rank}     & 0.396 & 0.094 & 0.302 & 0.328 & inherited \\
                & \stage{Explain}  & 0.225 & 0.067 & 0.158 & 0.008 & inherited \\
                & \stage{Memory}   & 0.156 & 0.056 & 0.101 & 0.000 & inherited \\
    \midrule
    Restaurants & \stage{Elicit}   & 0.056 & 0.056 & 0.000 & 0.000 & invariant \\
                & \stage{Retrieve} & 0.220 & 0.146 & 0.074 & 0.120 & direct \\
                & \stage{Rank}     & 0.490 & 0.156 & 0.333 & 0.370 & mixed \\
                & \stage{Explain}  & 0.197 & 0.072 & 0.125 & 0.014 & inherited \\
                & \stage{Memory}   & 0.160 & 0.056 & 0.104 & 0.004 & inherited \\
    \bottomrule
  \end{tabular}
\end{table}

\subsubsection{Localization}
\label{sec:scaff-localization}

Selecting, per scenario, the stage with the largest $\DA_s$---and \emph{none} when no stage
exceeds $\tau$---recovers the planted stage in all $48$ scenarios across the two domains,
including the control scenarios in which nothing was planted. This is the sanity check that has
to pass before the decomposition is applied to a real agent, where the causal origin is not
known.

Two qualifications keep the number honest. First, accuracy is reported at a single planted effect
size, and the planted faults here are strong enough to change the primary coordinate
substantially; localization in the weak-effect regime---the regime a deployed audit operates
in---is characterised by accuracy as a function of effect size, not by a single percentage, and we
report it that way in Section~\ref{sec:results}. Second, localization inherits the sensitivity of
whichever coordinate a stage uses. The top-one gap we adopt for \stage{Rank} is a binary
coordinate, and its noise floor in Table~\ref{tab:scaff-coords} ($0.328$, $0.370$) exceeds its own
direct effect: it is simultaneously noisy under repetition and insensitive to a genuine
reordering that leaves the head of the list intact. A graded coordinate---rank correlation, or a
difference in $\NDCG@K$---is the better choice for that stage, and the composition of
Table~\ref{tab:scaff-coords} at \stage{Rank} should be read with that limitation in view.

\subsubsection{Process Sensitivity Masked at the Endpoint}
\label{sec:scaff-masking}

Table~\ref{tab:scaff-masking} cross-tabulates whether a pair is process-sensitive---some stage
with $\DA_s > \tau$---against whether its final ranked list differs beyond the endpoint
equivalence margin. Of the $151$ process-sensitive pairs, $77$ are endpoint-equivalent: for just
over half of the trajectories in which some stage responded directly to the descriptor, the final
recommendation is indistinguishable. These are exactly the cases an output-level consumer-fairness
audit records as fair, and they are the empirical argument for auditing the process. Aggregated
per domain, $36.5\%$ (movies) and $43.8\%$ (restaurants) of all pairs are masked in this sense.

\begin{table}[t]
  \caption{Process sensitivity against endpoint difference, pooled across domains. The $77$ pairs
    in the upper right are process-sensitive trajectories that an output-only audit scores as
    fair.}
  \label{tab:scaff-masking}
  \begin{tabular}{lrr}
    \toprule
     & endpoint differs & endpoint equivalent \\
    \midrule
    process sensitive        & 74 & 77 \\
    no direct sensitivity    &  0 & 41 \\
    \bottomrule
  \end{tabular}
\end{table}

Part of this masking is structural rather than incidental, and the distinction matters for how
strongly the result can be stated. \stage{Explain} and \stage{Memory} both act downstream of the
ranked list, so a fault planted at either cannot alter the endpoint within a turn under any
parameterisation. For those stages an endpoint audit is not underpowered but blind by
construction; a \stage{Memory} fault becomes observable only once a later turn consumes the
corrupted profile. The empty lower-left cell is a soundness property of this substrate---a planted
fault is the only route by which the two arms can differ---and not a general guarantee. Under a
stochastic agent that cell will be populated, and it should be reported rather than assumed empty.

\subsubsection{Repair Removability}
\label{sec:scaff-repair}

Replacing the implicated stage's output with its valid counterpart and re-running all descendants
reduces the final ranking divergence by $0.576$ (movies) and $0.588$ (restaurants). The same
operation applied at a stage the crossover exonerates reduces it by $0.170$ and $0.187$. The
attribution therefore predicts the effect of acting on it, which a coordinate that merely
described the divergence would not.

The constraint noted in Section~\ref{sec:scaff} is visible in these results and is why we state
it: a repair at \stage{Rank} scores a removability exactly equal to its own pre-repair divergence
in both domains, because substituting the ranked list is substituting the quantity the outcome
metric is computed on. That figure measures the metric's definition rather than the repair's
effect, and removability for \stage{Rank} is accordingly evaluated on the following turn.

\subsubsection{Validity}
\label{sec:scaff-validity}

The validity gate rejected $2.1\%$ of \stage{Explain} outputs in both domains---all of them a
mismatch between the item named in the explanation and the top-ranked item, and all concentrated
on one scenario per domain. These runs are excluded from \stage{Explain}'s fairness coordinates
rather than scored, and the rejection rate is reported alongside them: a stage that frequently
emits structurally invalid output is a finding in its own right, and silently averaging over such
runs would report it as invariance.
